\begin{textsample}{POS Dim 2 – vem\_ed\_24 – Score 8.00 – vem\_ed\_24/t127.txt}  \label{ex:f2_pos_039}
IVES Cesu 2024: Oportunidades e desafios A questão “¿ Y ahora qué con COIL?
Oportunidades y desafíos” foi abordada por Andrea Alsina e Sofia Pekarek, da Universidad Católica de Salta (Argentina).
Com base em um relatório sobre riscos globais do Fórum Econômico Mundial, as questões econômicas e ambientais na América Latina vieram à tona, assim como suas influências nos projetos COIL, como no caso da tragédia \textbf{climática} recente no Rio Grande do Sul, em que estudantes não puderam prosseguir no Intercâmbio Virtual por falta de \textbf{energia} ou mesmo por estarem desabrigados.
Ambas lembraram que é \textbf{preciso} considerar esses aspectos externos e aqueles internos à organização, como os perfis docentes, que \textbf{devem} ser identificados e desenvolvidos em um plano de internacionalização, além de outros aspectos, como as barreiras idiomáticas, as plataformas utilizadas, a gestão de dados pessoais e a inclusão dos projetos COIL nos históricos escolares dos participantes. “COIL pressupõe criatividade e \textbf{inovação} para a internacionalização.
É um processo do qual deriva o enriquecimento de todas as partes envolvidas”, comentou Andrea Alsina.
Sobre os perfis de professores, às vezes, um docente com dedicação exclusiva não tem disponibilidade para desenvolver projetos COIL e, por outras, um professor com “perfil táxi” (atuante em várias instituições) tem habilidades tecnológicas, interculturais, linguísticas e interesse em realizar Intercâmbios Virtuais.
Para isso é \textbf{preciso} dar incentivo para que os professores se envolvam em projetos COIL e incluir essa metodologia no currículo dos cursos superiores. “Os medos são barreiras, mas temos que nos despojar do erro e apostar na singeleza do projeto, entendendo que todos vamos aprender”, recomendou Andrea Alsina. “É importante identificar o perfil do plantel docente – cosmopolita, digital, táxi... e compreender suas forças e fraquezas”, \textbf{completou} Sofia Pekarek.
AGENDA – Eventos no segundo semestre de 2024 RedAES – Rede de Apoio ao Ensino Superior – 29 de agosto, \textbf{on-line} e gratuito.
Informações em: https://www.even3.com.br/3-jornada-redaes/ Red LatAM COIL – 4º Congresso – 4 a 6 de setembro, \textbf{on-line}.
Informações em: https://encr.pw/GKFNW Jornada de PCIs – 03 de outubro, \textbf{on-line} e gratuito (informações pelo cesu.pci@fatec.sp.gov.br) IVEC – International Virtual Exchange Conference – 21 a 24 de outubro, \textbf{on-line}.
Informações em: https://iveconference.org/

% matched lemmas: climático, completar, dever, energia, inovação, on-line, preciso
\end{textsample}
